\documentclass[11pt]{article}

\usepackage[margin=1in]{geometry}
\usepackage[T1]{fontenc}
\usepackage[utf8]{inputenc}
\usepackage{booktabs}
\usepackage{array}
\usepackage{tabularx}
\usepackage{longtable}
\usepackage{enumitem}
\usepackage{hyperref}

\hypersetup{
  colorlinks=true,
  linkcolor=black,
  urlcolor=blue,
  pdftitle={Harness, Formalizer, and Agentic Prover Audit},
  pdfauthor={LeanEcon internal engineering note}
}

\newcolumntype{Y}{>{\raggedright\arraybackslash}X}

\title{Harness, Formalizer, and Agentic Prover Audit}
\author{LeanEcon internal engineering note}
\date{March 25, 2026}

\begin{document}
\maketitle

\section*{Executive Summary}

LeanEcon's current backend is not architecturally broken. The benchmark harness,
bounded formalizer, and agentic prover are separated cleanly enough that the
system can be hardened in place. The main problems are not missing modules;
they are missing context transfer, slow fallback paths, and expensive proving
loops whose search quality is still weak relative to their runtime.

The core diagnosis is straightforward:

\begin{enumerate}[leftmargin=1.5em]
  \item The formalizer is not receiving the full benefit of Lean-aware search.
  Its retrieval layer is mostly prompt-time only, optional MCP search is off by
  default in code, and the fast \texttt{lean\_run\_code} validation path is
  effectively unavailable in the benchmark evidence.
  \item The prover does real work, but it starts cold on every run and never
  inherits the formalizer's structured retrieval context. It receives a theorem
  stub, not the imports, identifiers, shape guidance, or diagnosis that led to
  that stub.
  \item The live Lovable integration adds a frontend handoff problem on top of
  the backend issues: explicit preamble choices are not preserved consistently,
  and the classify path can broaden a user-selected preamble into a noisier
  request payload.
\end{enumerate}

This document isolates the engineering system behind the harness, formalizer,
and prover; separates strengths from current limitations; explains the most
likely causes of runtime slowness; and ends with a prioritized remediation
stack.

\section{System Boundary}

\subsection*{Subsystem ownership}

\begin{tabularx}{\textwidth}{>{\raggedright\arraybackslash}p{0.18\textwidth} >{\raggedright\arraybackslash}p{0.34\textwidth} Y}
\toprule
Subsystem & Code ownership & Role in the system \\
\midrule
Benchmark harness & \texttt{src/benchmark\_harness.py} & Runs lane-based backend evaluations, aggregates latency, pass rates, stop reasons, partial runs, repair buckets, retrieval counts, semantic alignment, and provider telemetry. It measures backend behavior, not Lovable UI handoffs. \\
Formalizer & \texttt{src/formalizer.py}, \texttt{src/formalization\_search.py} & Cleans natural-language claims, builds bounded retrieval context, selects preambles, prompts Leanstral, applies deterministic repairs, and validates theorem stubs with \texttt{sorry}. \\
Agentic prover / verifier & \texttt{src/agentic\_prover.py}, \texttt{src/mcp\_runtime.py}, \texttt{src/lean\_verifier.py} & Drives tool-using proof search through Mistral Conversations plus \texttt{lean-lsp-mcp}, then performs the final truth check with local Lean compilation. \\
\bottomrule
\end{tabularx}

\subsection*{Runtime path from Lovable UI to Lean kernel}

\begin{enumerate}[leftmargin=1.5em]
  \item The Lovable frontend sends a claim to \texttt{/api/v1/classify} or
  skips directly to \texttt{/api/v1/formalize}.
  \item The formalizer builds a bounded retrieval context from keyword-ranked
  preambles, curated static hints, and optional MCP search hits.
  \item Leanstral produces a theorem stub ending in \texttt{sorry}; the
  formalizer may repair and revalidate it.
  \item The API queues \texttt{/api/v1/verify}, which calls
  \texttt{run\_pipeline(..., preformalized\_theorem=...)}.
  \item The agentic prover creates a fresh working Lean file, a fresh Mistral
  RunContext, and a fresh MCP client registration.
  \item The prover iterates between tactic writes and Lean diagnostics until it
  either reaches a proof, times out, or trips an internal budget guard.
  \item Final truth comes from \texttt{lake env lean} on an isolated per-run
  Lean file, not from MCP.
\end{enumerate}

\section{Strengths of the Current Design}

\begin{itemize}[leftmargin=1.5em]
  \item The module boundaries are good. Classification, formalization,
  retrieval, proving, verification, caching, and job orchestration are
  separated rather than tangled together.
  \item Final truth is anchored in local Lean compilation. MCP is advisory and
  interactive; it is not the trust root.
  \item Verification uses unique temporary Lean files rather than a shared
  tracked module, which makes concurrent verify jobs far safer than the old
  single-file path.
  \item The harness is already more honest than many systems at this stage. It
  splits lanes, records partial outcomes, exposes repair buckets, and tracks
  provider telemetry rather than hiding weak paths behind one blended metric.
  \item The prover is genuinely MCP-backed. This is not paper architecture:
  tool traces, diagnostic checks, search budgets, and interruption handling all
  exist in the implementation.
  \item The preamble library, catalog generation, and bounded retrieval helpers
  give the system a reusable domain substrate instead of asking the model to
  reinvent all economics definitions from scratch every time.
\end{itemize}

\section{Limitations of the Current Design}

\begin{itemize}[leftmargin=1.5em]
  \item Formalization retrieval is shallow. The system can inject imports and
  prompt hints, but it does not yet use retrieval as an active search-and-check
  loop.
  \item The formalizer and prover are loosely coupled. The pipeline hands the
  prover only theorem text, so the prover must rediscover context the
  formalizer already assembled.
  \item Public API models strip away the deep telemetry needed to debug why a
  given claim was slow or why retrieval was skipped. The code keeps that
  telemetry internally, but the frontend cannot see it.
  \item Preamble matching uses two different policies: a broad classifier-facing
  matcher and a tighter formalizer auto-selection matcher. That mismatch creates
  avoidable noise.
  \item The benchmark harness measures backend lanes, but it does not cover the
  Lovable UI branch logic that decides whether explicit preambles are preserved,
  widened, or dropped before \texttt{/formalize}.
  \item Agentic proof search still buys too little for what it costs on harder
  claims. The current trace data shows many tool cycles ending in no verified
  progress.
\end{itemize}

\section{Why the System Is Slow}

\subsection*{1. Fast formalizer validation is not carrying the load}

The intended fast path for \texttt{sorry}-validation is
\texttt{lean\_run\_code} through MCP. In practice, the benchmark evidence says
that the formalizer is almost always falling back to local
\texttt{lake env lean} checks instead.

The latest tier-1 formalizer-only report shows:

\begin{itemize}[leftmargin=1.5em]
  \item \texttt{validation\_method = lake\_env\_lean} on all 6 attempts
  \item \texttt{validation\_fallback\_reason = lean\_run\_code\_unavailable} on all 6 attempts
  \item a p50 formalizer latency of 18{,}431.8 ms and a p95 of 80{,}031.8 ms
\end{itemize}

This is consistent with the code comments in \texttt{src/lean\_verifier.py},
the tests in \texttt{tests/test\_lean\_runner.py}, and multiple snapshot
artifacts that record MCP-side failure text of the form ``No valid Lean project
path found. Run another tool first to set it up.'' The result is that the
formalizer is paying a full local compiler check far more often than intended.

\subsection*{2. Formalization MCP retrieval is off by default}

The retrieval helper in \texttt{src/formalization\_search.py} sets
\texttt{LEANECON\_ENABLE\_FORMALIZATION\_MCP\_SEARCH} to false unless the
environment overrides it. During this audit, the repo-local environment files
did not define a formalization retrieval flag, so the code default is the
effective default. Even when enabled, the search budget is narrow:

\begin{itemize}[leftmargin=1.5em]
  \item at most 2 MCP search queries
  \item a 5-second overall retrieval timeout
  \item prompt-time hits only, not a persistent retrieval state carried into proving
\end{itemize}

So the system's ``Mathlib search engine'' is present in code but not operating
as a first-class runtime assistant.

\subsection*{3. Every proving job starts cold}

The prover deliberately creates a fresh RunContext and a fresh MCP client for
each run because the current SDK lifecycle does not safely support warm reuse.
That is a pragmatic choice, but it means every verify job pays:

\begin{itemize}[leftmargin=1.5em]
  \item working file initialization
  \item MCP tool registration
  \item MCP session startup and priming
  \item a brand new Conversations loop
\end{itemize}

This startup cost matters because proving is already the slowest lane even when
formalization is bypassed. In the latest tier-1 full benchmark report:

\begin{itemize}[leftmargin=1.5em]
  \item \texttt{theorem\_stub -> verify} had p50 174{,}297.9 ms and p95 180{,}402.5 ms
  \item \texttt{raw\_lean -> verify} had p50 138{,}305.2 ms and p95 146{,}799.8 ms
\end{itemize}

Those numbers show that prover-plus-verifier latency is a major bottleneck even
when the claim-shaping stage is skipped.

\subsection*{4. \texttt{TIMEOUT\_MS} is not a true total proof wall-clock cap}

This is a code-based inference, but it is an important one. In
\texttt{src/agentic\_prover.py}, \texttt{TIMEOUT\_MS = 120\_000} is commented as
``2 minutes for the full run\_async conversation.'' The implementation does not
enforce it that way. Instead, the value is passed as \texttt{timeout\_ms} to
each individual \texttt{conversations.start\_async} and
\texttt{conversations.append\_async} request.

There is no outer \texttt{asyncio.wait\_for(...)} or other wall-clock guard
around the entire proving session. That means aggregate proof runtime can
substantially exceed 120 seconds if the loop keeps making successful append
calls. The benchmark latencies above are consistent with that reading.

\subsection*{5. The prover spends too many append cycles for the progress it gets}

The local trace summary from \texttt{logs/runs.jsonl} is stark:

\begin{itemize}[leftmargin=1.5em]
  \item 86 runs parsed
  \item 959 tool calls
  \item 83 successful tactic applications
  \item tool-call efficiency 0.087
  \item tool-call waste ratio 0.913
\end{itemize}

The tool mix is also revealing:

\begin{itemize}[leftmargin=1.5em]
  \item 293 writes
  \item 312 diagnostics
  \item 70 search
  \item 40 blocked calls
\end{itemize}

This is not a picture of a search loop that is quickly converging on proofs. It
is a picture of a loop that is spending many expensive interaction cycles to get
very little verified proof progress.

\subsection*{6. Interrupted runs still pay for final verification}

On timeout and other interruption paths, the prover calls back into
\texttt{verify(current\_code)} on the partial state before returning. This is
useful for salvaging partial value, but it means the system can pay for a final
Lean compilation even after the proving loop has already gone long.

\section{Why It Feels Disconnected from Preamble and Mathlib Search}

\subsection*{1. Retrieval is present, but mostly as prompt text}

The current formalization retrieval stack consists of:

\begin{itemize}[leftmargin=1.5em]
  \item keyword-ranked preambles from \texttt{src/preamble\_library.py}
  \item curated static import and identifier hints from \texttt{src/formalization\_search.py}
  \item optional MCP \texttt{lean\_local\_search} hits when retrieval is enabled
\end{itemize}

That context is inserted into the prompt, and preamble imports are injected into
the theorem stub. What it does not become is a persistent structured object that
the prover can continue using.

\subsection*{2. The prover never receives the formalizer's retrieval state}

The key pipeline break is that \texttt{prove\_and\_verify()} receives only
\texttt{theorem\_with\_sorry}. It does not receive:

\begin{itemize}[leftmargin=1.5em]
  \item selected preamble names
  \item candidate imports
  \item candidate identifiers
  \item shape guidance
  \item retrieval notes
  \item formalization diagnosis or repair history
\end{itemize}

As a result, the prover starts from a string and a working file path. It has to
re-infer proof structure from the generated theorem rather than inheriting the
formalizer's rationale.

\subsection*{3. The prover's search surface is much narrower than the formalizer's ideal search surface}

The agentic prover prunes its tools down to:

\begin{itemize}[leftmargin=1.5em]
  \item \texttt{apply\_tactic}
  \item \texttt{lean\_goal}
  \item \texttt{lean\_diagnostic\_messages}
  \item \texttt{lean\_multi\_attempt}
  \item \texttt{lean\_code\_actions}
  \item \texttt{lean\_state\_search}
  \item \texttt{lean\_hammer\_premise}
\end{itemize}

Notably absent are the broader discovery tools that would make the system feel
connected to a true Mathlib search workflow, such as \texttt{lean\_local\_search},
completions, finder-style semantic search, or import discovery.

\subsection*{4. The preamble policies are mismatched}

This is one of the clearest design disconnects in the current system:

\begin{itemize}[leftmargin=1.5em]
  \item \texttt{classify\_claim()} uses \texttt{find\_matching\_preambles()},
  which is based on the broad \texttt{keywords} list.
  \item \texttt{build\_formalization\_context()} uses
  \texttt{rank\_matching\_preambles(auto=True)}, which uses tighter
  \texttt{auto\_keywords} and caps automatic selection at 2 entries.
\end{itemize}

For CRRA-style claims, the broad classifier matcher can pull in
\texttt{cara\_utility} because the keywords overlap on phrases such as
``risk aversion'' and ``derivative.'' The formalizer's auto-selection path is
tighter and, by test design, tries to avoid that noise.

\subsection*{5. The public API hides the telemetry needed to debug retrieval}

\texttt{FormalizeResponse} exposes \texttt{provider\_telemetry} but not the
full \texttt{formalizer\_telemetry} object. \texttt{VerifyResponse} likewise
omits deep trace fields such as \texttt{tool\_trace}, \texttt{tactic\_calls},
and the agent summary that exist internally.

This matters because a frontend engineer looking only at the public API cannot
see:

\begin{itemize}[leftmargin=1.5em]
  \item whether MCP retrieval ran
  \item why retrieval was skipped
  \item which validation method was used
  \item how many repair buckets fired
  \item how much proof search was blocked or wasted
\end{itemize}

\subsection*{6. The Lovable handoff really is inconsistent}

Live Playwright reproduction on \url{https://leanecon-demo.lovable.app/} on
March 25, 2026 showed three different handoff behaviors:

\begin{enumerate}[leftmargin=1.5em]
  \item \textbf{Theorem-card \texttt{Pipeline}} initially populated the claim
  text but did not visibly preload preamble context.
  \item \textbf{\texttt{Use 1 in Pipeline}} visibly loaded exactly one explicit
  preamble (\texttt{crra\_utility}) before the run started.
  \item After clicking the normal classify path from that same one-preamble
  starting point, the frontend sent \texttt{/formalize} a broadened payload of
  three explicit preambles:
  \texttt{crra\_utility}, \texttt{arrow\_pratt\_rra}, and \texttt{cara\_utility}.
  \item After clicking \texttt{Skip classify} from both theorem-card and
  one-preamble starting points, the frontend sent \texttt{/formalize} no
  explicit \texttt{preamble\_names} at all and relied on backend auto-selection.
\end{enumerate}

That is not just a cosmetic UI inconsistency. It changes the actual backend
request body and therefore changes prompt context and theorem shape.

\section{Evidence}

\subsection*{Benchmark lane summary}

\begin{table}[htbp]
\centering
\small
\begin{tabularx}{\textwidth}{>{\raggedright\arraybackslash}p{0.34\textwidth} >{\raggedleft\arraybackslash}p{0.12\textwidth} >{\raggedleft\arraybackslash}p{0.12\textwidth} >{\raggedleft\arraybackslash}p{0.12\textwidth} Y}
\toprule
Source & pass@1 & p50 ms & p95 ms & Main signal \\
\midrule
Tier-1 full: raw\_claim -> full API & 0.333 & 367266.9 & 428133.6 & Weakest lane; claim shaping plus proof search is still the unstable path. \\
Tier-1 full: theorem\_stub -> verify & 1.000 & 174297.9 & 180402.5 & Even strong prove-only runs are very slow, so prover/verifier cost dominates. \\
Tier-1 full: raw\_lean -> verify & 1.000 & 138305.2 & 146799.8 & Bypassing formalization helps, but verify remains expensive. \\
Tier-1 formalizer-only gate & 0.833 & 18431.8 & 80031.8 & Formalization is improving, but still falls back to slow local validation. \\
\bottomrule
\end{tabularx}
\caption{Latest benchmark evidence used in this audit.}
\end{table}

\subsection*{Trace summary from \texttt{logs/runs.jsonl}}

\begin{table}[htbp]
\centering
\small
\begin{tabularx}{\textwidth}{>{\raggedright\arraybackslash}p{0.42\textwidth} >{\raggedleft\arraybackslash}p{0.18\textwidth} Y}
\toprule
Metric & Value & Interpretation \\
\midrule
Runs parsed & 86 & Sufficient to show proving behavior is not a one-off anecdote. \\
Total tool calls & 959 & The prover is doing a large amount of interactive work. \\
Successful tactic applications & 83 & Only a small fraction of tool activity ends in successful tactic steps. \\
Tool-call efficiency & 0.087 & The search loop is low-yield. \\
Tool-call waste ratio & 0.913 & Most tool activity does not translate into successful tactic progress. \\
Tool mix & 293 writes, 312 diagnostics, 70 search, 40 blocked & Diagnostics and loop control consume a large share of runtime. \\
\bottomrule
\end{tabularx}
\caption{Trace evidence from \texttt{scripts/analyze\_traces.py --runs-file logs/runs.jsonl --format text}.}
\end{table}

\subsection*{Live Lovable reproduction}

\begin{table}[htbp]
\centering
\small
\begin{tabularx}{\textwidth}{>{\raggedright\arraybackslash}p{0.29\textwidth} >{\raggedright\arraybackslash}p{0.33\textwidth} Y}
\toprule
Observation & Concrete evidence & Implication \\
\midrule
One explicit preamble was not stable through classify & \texttt{Use 1 in Pipeline} loaded \texttt{crra\_utility}, but the subsequent \texttt{/formalize} request sent three explicit preambles. & Classify-path request assembly can widen user intent and inject noise. \\
Skip classify dropped explicit preambles entirely & \texttt{/formalize} request body contained only \texttt{raw\_claim}, even when one preamble had been preloaded in the UI. & The explicit preamble handoff is broken on this branch. \\
Proof search timed out on a mid-level CRRA case & Verify ended with \texttt{partial = true}, \texttt{stop\_reason = timeout}, \texttt{phase = proved}. & Timeouts are not limited to frontier claims. \\
The timeout was expensive before failing & Live CRRA run used 37 Conversations API calls and about 42.2 s of provider latency inside a 75.9 s job. & A large fraction of wall-clock time is consumed before the final Lean rejection. \\
\bottomrule
\end{tabularx}
\caption{Live observations from \url{https://leanecon-demo.lovable.app/} on March 25, 2026.}
\end{table}

\subsection*{Worked example: CRRA reproduction}

The clearest live reproduction from this audit used the claim ``Under CRRA
utility, relative risk aversion equals gamma.'' The request path and outcome
were:

\begin{enumerate}[leftmargin=1.5em]
  \item The classify path forwarded three explicit preambles:
  \texttt{crra\_utility}, \texttt{arrow\_pratt\_rra}, and
  \texttt{cara\_utility}.
  \item The verify job ended with \texttt{partial = true},
  \texttt{stop\_reason = timeout}, \texttt{phase = proved}, and
  \texttt{proof\_generated = true}.
  \item The final kernel error was
  \texttt{field\_simp made no progress on goal}.
  \item The partial proof still contained a nested \texttt{sorry}, and the UI
  warning reported that the prover stopped during Lean/MCP cleanup after a timeout.
  \item Provider telemetry for that run reported 37 Conversations API calls:
  1 start plus 36 append calls.
\end{enumerate}

This example is useful because it was not an exotic frontier theorem. It was a
preamble-backed, mid-level economics claim that should have benefited from the
system's domain substrate, yet the runtime still degraded into an expensive
partial proof.

\section{Strategy}

The highest-ROI remediation stack is:

\begin{enumerate}[leftmargin=1.5em]
  \item \textbf{Restore a reliable fast formalizer validation path.} Make
  \texttt{lean\_run\_code} dependable again or replace it with an equally fast
  project-aware validator. Until that is fixed, the formalizer will keep paying
  local compiler costs too early and too often.
  \item \textbf{Make formalization retrieval truly active and measurable.}
  Turn MCP retrieval on deliberately, expand beyond static hints when useful,
  and expose whether it ran, what it returned, and why it was skipped.
  \item \textbf{Pass structured formalizer context into the prover.} The prover
  should inherit selected preambles, candidate identifiers, theorem-shape
  guidance, and repair diagnosis rather than re-deriving them from raw text.
  \item \textbf{Fix frontend preamble handoff and merge policy.} A user-selected
  preamble should either be preserved exactly or deliberately merged according
  to a documented policy. It should not silently disappear on one branch and
  silently widen on another.
  \item \textbf{Add a real proof wall-clock budget and an append-round cap.}
  Keep per-request API timeouts, but add an outer cap on total proving time and
  total model round-trips.
  \item \textbf{Expose debugging telemetry in the public API and UI.} At
  minimum, surface validation method, retrieval usage, stop reason, partial
  proof metadata, and a compact trace summary.
  \item \textbf{Tighten benchmark and trace regression gates.} Add explicit
  regression targets for timeout frequency, partial proof rate, append-round
  counts, and tool-call efficiency on preamble-backed claims.
\end{enumerate}

\section*{Bottom Line}

LeanEcon's engineering foundation is good enough to improve in place. The
system is slow and feels disconnected from its search and preamble assets not
because those assets are absent, but because they are only partially activated,
not carried across subsystem boundaries, and not surfaced clearly enough to the
frontend. The path forward is to make retrieval real, make context persistent,
and make proving budgets honest.

\end{document}
